% ===== Chapter 6: auto-converted from Word =====

\section{Conclusion}

This thesis asked whether aggregate insider ownership predicts risk-adjusted stock returns in a global sample. Wood (2025) had shown that owner operated firms compound faster than the global benchmark on a raw return basis. The question of whether that outperformance survives risk adjustment remained open. The five factor portfolio sort approach applied here closes that question. The universe earns positive alpha at every ownership level and the within universe concentration premium is statistically and economically meaningful.

The importance of this result is that the alignment mechanism operates at global scale and is detectable in market prices. Four decades of governance research have discussed whether concentrated insider ownership benefits or harms outside shareholders. The theory split in the literature between convergence of interest and entrenchment predictions. The evidence here weighs in favour of convergence: the highest ownership decile earns the largest risk-adjusted return, with no clear sign of an entrenchment turning point. This challenges the assumption built into many institutional governance papers, that majority insider control is unfavourable to minority shareholders. This provides an empirical basis for treating concentrated insider equity as a positive rather than negative signal in investment screening.

What follows from this work is a set of questions which are unanswered. Aggregate insider ownership now appears as a firm characteristic that is relevant for returns. This may support the need for treatment alongside size, value, and profitability in factor model construction. Whether ownership is a priced systematic risk or a proxy for the high ownership characteristics discussed above is not yet identified. This remains an open question that an constructed global ownership factor could address. A specification using changes in ownership would test whether the information content comes from the ownership change and if the relationship is causal. A regression design modelling dynamic ownership would identify the causal channel that the portfolio sort approach used here cannot. Each future paper extends directly from what this thesis establishes and brings the global ownership return relationship closer to the causal understanding that the underlying corporate governance theories ultimately need. 
